\documentclass[x11names,svgnames]{beamer}

\usetheme{metropolis}
\setbeamertemplate{caption}{\raggedright\insertcaption\par}
\setbeamercolor{block title}{fg=LightSkyBlue4}
\newenvironment<>{importantblock}[1]{%
  \begin{actionenv}#2%
    \def\insertblocktitle{#1}%
    \par%
    \mode<presentation>{%
      \setbeamercolor{block title}{%
        use=normal text,
        fg=normal text.fg,
        bg=normal text.bg!80!fg
      }
      \setbeamercolor{block body}{
        use={block title, normal text},
        bg=block title.bg!50!normal text.bg
      }
    }%
    \usebeamertemplate{block begin}}
  {\par\usebeamertemplate{block end}\end{actionenv}}
\newenvironment<>{decisionblock}[1][Statement of the problem]{%
  \begin{actionenv}#2%
    \def\insertblocktitle{#1}%
    \par%
    \mode<presentation>{%
      \setbeamercolor{block title}{%
        use=normal text,
        fg=normal text.fg,
        bg=normal text.bg!80!fg
      }
      \setbeamercolor{block body}{
        use={block title, normal text},
        bg=block title.bg!50!normal text.bg
      }
    }%
    \usebeamertemplate{block begin}}
  {\par\usebeamertemplate{block end}\end{actionenv}}
\usepackage{fontspec}
\usepackage[french, frenchkw]{algorithm2e}
\SetKwFor{Pc}{Pour chaque}{faire}{fin}
\SetKwProg{Fn}{Fonction}{:}{}
\usepackage{textcomp}
\usepackage{minted}
\usepackage{pdfpages}
\usepackage{qrcode}
\usepackage[francais]{babel}
\usepackage{xspace}
\usepackage[normalem]{ulem}
\usepackage[backend=biber,style=authoryear,url=false]{biblatex}
\addbibresource{Bibliography.bib} % Bibliography file
\usepackage{amsmath}
\usepackage{mathtools}
\usepackage{pifont}% http://ctan.org/pkg/pifont
\newcommand{\cmark}{\ding{51}}%
\newcommand{\xmark}{\ding{55}}%
\usepackage{marvosym}
\usepackage{tikz}
\usetikzlibrary{calc, positioning, arrows, automata, fadings, decorations.pathmorphing, shapes, patterns, tikzmark, fit, chains, snakes, shapes.multipart}
% https://tex.stackexchange.com/questions/401884/how-do-i-change-hyperlinks-color-only
\hypersetup{
  colorlinks,
  allcolors=.,
  urlcolor=DarkBlue,
}


\author{}
\date{Lundi 25 septembre 2023}

\title{\includegraphics[width=0.75\textwidth]{../../img/logo-igm.png} \\ Algorithmique et Programmation 1}
\institute{L1 Mathématiques - L1 Informatique \\ Semestre 1}

\begin{document}

\maketitle

\section{Chapitre 4 : Fonctions}

\begin{frame}{Conclusions de l'activité}
  \begin{itemize}
  \item Chaque chef de rang a exécuté la même consigne
  \item Qui dépend possiblement du numéro de rangée
  \end{itemize}

  En pseudo-code, ça pourrait donner :

  \begin{algorithm}[H]
    \Pc{rangée $r$}{
      \emph{dessiner\textunderscore ligne}(r)\;
    }
  \end{algorithm}

  \begin{algorithm}[H]
    \Fn{dessiner\textunderscore ligne(r)}{
    i $\leftarrow$ 0\;
    \Tq{$i < n$}{
      Faire lever une feuille rose\;
      Avancer d'un.e étudiant.e\;
      i $\leftarrow i + 1$
    }
    }
  \end{algorithm}
  
\end{frame}

\begin{frame}
  \begin{block}{Fonction}
    \vspace{0pt}
    En informatique, une fonction est :
    \begin{itemize}
    \item Un morceau de programme
    \item Portant en général \emph{un nom}
    \item Prenant un ou plusieurs \emph{paramètres} (ou zéro)
    \item Renvoyant un résultat (la plupart du temps)
    \end{itemize}
  \end{block}
  \MVRightarrow{} Ça ressemble beaucoup aux fonctions vues en mathématiques !
\end{frame}

\begin{frame}
  \frametitle{À quoi ça sert ?}

  \begin{block}{Lisibilité}
    \begin{itemize}
    \item Isoler les parties du programme selon ce qu'elles font
    \item Éviter trop d'imbrications entre \mintinline{python}{if} et \mintinline{python}{while}
    \end{itemize}
  \end{block}

  \pause

  \begin{block}{Modularité et robustesse}
    \begin{itemize}
    \item Réutiliser le code plusieurs fois
    \item Faciliter la correction des bugs
    \end{itemize}
  \end{block}

  \pause
  \begin{block}{Généricité}
    \begin{itemize}
    \item Changer la valeur des paramètres
    \end{itemize}
  \end{block}
\end{frame}

\begin{frame}{Fonctions prédéfinies}
  \MVRightarrow{} On en a déjà vu !

  \begin{exampleblock}{Exemples}
    \begin{itemize}
    \pause
  \item \mintinline{python}{print(s)} : afficher la chaîne de caractères \mintinline{python}{s}
  \item \mintinline{python}{input(s)} : afficher la chaîne de caractères \mintinline{python}{s} et attendre une entrée utilisateur
  \item \mintinline{python}{int(x)} : convertit l'objet \mintinline{python}{x} en entier
  \item \mintinline{python}{len(x)} : renvoie la longueur de l'objet \mintinline{python}{x}
  \end{itemize}
\end{exampleblock}

\pause

\begin{block}{Modules}
  \vspace{0pt}
  Bibliothèque de fonctions, de types et d'objets
\end{block}

\begin{exampleblock}{Exemples}
  \pause
  \begin{itemize}
  \item \mintinline{python}{randint(a,b)} (module \mintinline{python}{randint}) : renvoie un entier aléatoire compris entre \mintinline{python}{a} et \mintinline{python}{b}
  \end{itemize}
  \end{exampleblock}
  
\end{frame}

\begin{frame}[fragile]
  \frametitle{Définition et appel de fonction}

  \begin{block}{Définir une fonction}
    \vspace{0pt}
    La syntaxe pour définir une fonction est :
    
\begin{minted}{python}
def nom_fonction(parametre_1, ..., parametre_n):
    # corps de la fonction
    # utilisant parametre_1, ..., parametre_n
    ...
    # peut renvoyer un résultat :
    return resultat
\end{minted}
  \end{block}
  \pause
  \begin{block}{Appeler une fonction}
    \vspace{0pt}
    On peut ensuite \emph{appeler} la fonction \mintinline{python}{nom_fonction} das le code :
\begin{minted}{python}
resultat = nom_fonction(expression_1, ..., expression_n)
\end{minted}
  \end{block}
\end{frame}

\section{Exemples}

\begin{frame}[fragile]{Fonction à paramètres et résultats}
  Calculer le maximum de deux entiers :
  \pause
\begin{minted}{python}
def maximum(a,b):
    if a >= b:
        return a
    else:
        return b
\end{minted}
\pause
  \begin{block}{Entraînement}
    \begin{itemize}
    \item Ouvrir le programme \texttt{maximum.py}
    \item Décrire l'exécution pas à pas du programme (avec état de la mémoire)
    \item On peut aussi essayer avec \href{https://pythontutor.com/}{Python Tutor}
    \item Dresser un tableau de valeurs de l'exécution du programme
    \end{itemize}
  \end{block}
\end{frame}

\begin{frame}[fragile]{Fonction sans paramètres}
  Simuler un lancé de dés, et compter combien de coups il faut pour faire un 6
  \pause
\begin{minted}{python}
from random import randint
def lance_de():
    return randint(1,6)

compteur = 1
while lance_de() != 6:
    compteur = compteur + 1
print('Obtenu un 6 en', compteur, 'jets de dé.')
\end{minted}
  \pause
  Ce que ça donne sur \href{https://pythontutor.com/}{Python Tutor}.
\end{frame}

\begin{frame}[fragile]{Fonction sans valeur de retour}
  Dessiner un carré fait du caractère \mintinline{python}{caractere} :
  \pause
\begin{minted}{python}
def dessine_carre(n, caractere):
    i = 0
    while i < n:
        j = 0
        while j < n:
            print(caractere, end = '')
            j += 1
        print('\n', end = '')
        i += 1
    return
\end{minted}
  \pause
Ce que ça donne sur \href{https://pythontutor.com/}{Python Tutor}.
\end{frame}
\begin{frame}[fragile]{Erreur fréquente : confusion \begin{tabular}{l}paramètre / saisie \\ retour / affichage\end{tabular}}

  \begin{block}{Paramètre / saisie}
\begin{minted}{python}
def maximum(a,b):
    a = int(input()) # NON !
    b = int(input()) # NON !
    if a >= b:
        return a
    else:
        return b
\end{minted}
  \end{block}

  \begin{block}{Retour / affichage}
\begin{minted}{python}
def maximum(a,b):
    if a >= b:
        print(a) # NON !
    else:
        print(b) # NON !
\end{minted}
  \end{block}
\end{frame}

\begin{frame}
  \frametitle{Composition de fonctions}
  
  On peut appeler une fonction dans une fonction ! (et ainsi de suite)
\end{frame}

\begin{frame}[fragile]{Dessiner un carré : variante}
\begin{minted}{python}
def dessine_ligne(n, caractere):
    j = 0
    while j < n:
        print(caractere, end = '')
        j += 1
    print('\n', end = '')
    return

def dessine_carre(n, caractere):
    i = 0
    while i < n:
        dessine_ligne(n,caractere)
        i += 1
    return
\end{minted}
\end{frame}

\begin{frame}
  \frametitle{Fonctions et espaces de nom}
  \begin{itemize}
  \item les paramètres et variables définies dans le corps d'une fonction sont indépendantes des autres variables du programme
  \item elles n'existent plus une fois l'exécution de la fonction terminée
  \item on les appelle des variables \emph{locales}
\end{itemize}

Par conséquent, on peut renommer les variables d'une fonction.

$\rightarrow{}$ Démonstration sur Thonny (\texttt{maximum.py}).

\end{frame}

\begin{frame}
  \frametitle{(Contre-)exemple : intervertir des variables}
  \begin{itemize}
  \item changer les valeurs de `a` et `b` dans la fonction n'a pas d'effet sur `x` et `y` dans le programme principal !
    \item la variable `temp` n'existe plus après l'exécution de la fonction
  \end{itemize}
\end{frame}
\end{document}

%%% Local Variables:
%%% TeX-engine: xetex
%%% TeX-command-extra-options: "-shell-escape"
%%% End: